\section{The full logical and arithmetic bridge}\label{sec:bridge}
\subsection{All intuitionistic operations}\label{sec:choice-introduction}
The intended comparison is with full first-order iPRA. The following
macros are the inherited induction-independent case-logic construction
\cite{choice}; they now apply because equality was already recovered.

Replace Q inside macros by its interderivable D form with unused witness.
On Q-free formulas define
\[
 G_v(s=t):=(C(v,0,s)=C(v,0,t)),\qquad
 G_v(\D_{x,y}B):=\D_{x,y}G_vB.
\]
With capture avoidance, computation, Ref and case introduction give
\[
 \vdash G_0A,\qquad G_{Sv}A\dashv\vdash A,\qquad
 A\dashv\vdash\D_{w,y}G_yA.
\]
For reverse padding, open D, supply its true zero row, and specialize at 1.

The selector between formulas is
\[
 M_v(s=t,r=q):=(C(v,s,r)=C(v,t,q))
\]
and, with aligned fresh binders,
\[
\begin{aligned}
 M_v(\D_{x,y}A,\D_{x,y}B)&:=\D_{x,y}M_v(A,B),\\
 M_v(E,\D_{x,y}B)&:=\D_{x,y}M_v(G_yE,B),\\
 M_v(\D_{x,y}A,E)&:=\D_{x,y}M_v(A,G_yE),
\end{aligned}
\]
where $E$ is an equation. Maximum input binder depth decreases.
Padding and binder congruence yield
$M_0(A,B)\dashv\vdash A$ and $M_{Sv}(A,B)\dashv\vdash B$.

Put $\delta(v)=C(v,1,0)$ and $p(z)=I(z,0;u.\delta(u))$;
$p(1)=1$ and $p(2)=0$ are computations. With fresh coordinates and unused
$w$, set
\[
\begin{aligned}
 \bot_*&:=(0=1),\\
 A\Rightarrow_*B&:=\D_{w,z}M_z(A,B),\\
 A\wedge_*B&:=\D_{w,z}G_zM_{p(z)}(A,B),\\
 \exists_*x\,A(x)&:=\D_{x,y}G_yA(x),\\
 A\vee_*B&:=\exists_*v\,M_v(A,B),\\
 \forall_*x\,A(x)&:=A(0)\wedge_*\D_{w,z}G_zA(z).
\end{aligned}                                                   \tag{Logic}
\]

Implication introduction is DI$_c$ using its zero row as the assumption;
elimination opens and specializes at 1. Existential introduction supplies
the witness, and elimination opens it freshly and removes the gate.
Disjunction injections use 0 and 1; elimination uses numerical cases on
the selector and the two retained branch implications.
Conjunction introduction uses selector cases and elimination uses rows
1 and 2. The universal keeps zero as a conjunct and supplies successors
in the other conjunct; coverage yields every instance.
Its introduction uses zero and successor substitutions of a proof at a
fresh parameter.

Coverage $(n=0)\vee_*\exists_*m(n=Sm)$ follows by cases and Ref.
From $0=1$, congruence in $v\mapsto C(v,s,t)$ gives every equation;
formula-structural induction gives explosion. These supply all ordinary
intuitionistic rules with the usual eigenvariable conditions.
No excluded middle for arbitrary formulas is added. The ordinary
realization has the explicit falsity bridge $(0=S0)\to\bot$.

\subsection{Scalar uniqueness and the noncircular induction bootstrap}
In U choose $a=f(0)$, use Ref for the base and reverse its conclusion:
\[
 \frac{\Gamma\vdash f(Sz)=h(f(z))}
      {\Gamma\vdash f(n)=I(n,f(0);u.h(u))}.                \tag{UR}\label{eq:ur}
\]
This is exactly the existing scalar uniqueness principle \cite{ur}.
Its equality and logical background has already been recovered.

Appendix~\ref{app:ur} gives the complete inherited proof that UR yields
equation induction. Its dependency order is important:
\begin{enumerate}
\item UR gives $I(n,0;u.Su)=n$ and retraction rigidity. If
$e(0)=0$, $e(e(x))=e(x)$ and $e(S(e(x)))=S(e(x))$, take
$h(u)=e(S(e(u)))$, $r(n)=I(n,0;h)$. Common-update uniqueness gives
$e(r(n))=r(n)$; then $r$ has successor update, so $r(n)=n$ and $e(n)=n$.
For $q(n)=0$, coverage and separation decide the property and the term
$e(n)=C(q(n),n,0)$ provides the retraction split. This derives contextual
induction for zero equations.
\item A scalar predecessor is constructed from finite clocks, a ramp, a
comparison iterator and bounded growth. Common-update uniqueness proves
$P0=0$ and $P(Sn)=n$, without using pairing, a general recursion compiler
or equation induction.
\item Define $d(a,b)=I(b,a;u.Pu)$ and
$\operatorname{sg}(a)=C(a,0,1)$. A scalar diagonal-descent argument and the
explicit accumulator
\[
 \mu(n,a)=d(n,d(n,a)),\quad
 \phi(n,a,b)=\mu(n,a)+\mu(d(n,a),d(b,a))
\]
prove $a+d(b,a)=b+d(a,b)$. Thus
\[
 \Delta(a,b)=\operatorname{sg}(d(a,b)+d(b,a))
\]
is Boolean and $\Delta(a,b)=0\leftrightarrow a=b$ on the entire domain.
\item Translate an equation-induction base and step for $s(n)=t(n)$ to
$\Delta(s(n),t(n))=0$, use zero-equation induction, and reflect back.
The retained context may contain arbitrary first-order formulas.
\end{enumerate}
The substantial scalar construction is inherited and attributed, not
claimed as a new result of U. Its full equations and proofs appear in
the appendix. Once equation induction is recovered, from
$\Gamma,E(z)\vdash E(Sz)$ work under $E(0)$ to derive $E(Sw)$ and use
conditional Q introduction. This recovers the old equation-progress rule.

\section{Equiconsistency and both interpretation composites}
\begin{theorem}[Exact native equivalence]
For every finite sequent of the repository's native calculus $\mathcal N$,
\[
 \Gamma\vdash_{\mathcal N}A
 \quad\Longleftrightarrow\quad
 \nf(\Gamma)\vdash_{\mathcal U}\nf(A).
\]
\end{theorem}
\begin{proof}
Translate native rules to normal forms. Computation becomes Ref,
transport is Section~\ref{sec:equality}, progress is
Section~\ref{sec:bridge}, and the other rules are unchanged.

Conversely, the native calculus proves normalization equalities and
derives U by equation induction on $I(z,a;h)=f(z)$.
Its base uses $f(0)=a$ and symmetry. Its successor step uses the
induction equation, congruence under $h$, computation and the retained
recurrence $f(Sz)=h(f(z))$. All the remaining U-calculus rules are native
rules or derived reflexivity.
\end{proof}

\begin{theorem}[Full first-order iPRA]
The ordinary realization of $\mathcal U$ is equiconsistent and
bi-interpretable with full first-order iPRA. Both numerical domains are
whole, equality is actual, and the two numerical composite isomorphisms
have graph $x=y$.
\end{theorem}
\begin{proof}
Compose the exact native equivalence with the scalar bridge
\cite{bridge}. In iPRA interpret each fixed iterator by primitive
recursion, and Q,D by \eqref{eq:read}. Equation induction validates U.

Conversely, predecessor and equality reflection supply successor
injectivity and characteristics for every fixed quantifier-free formula.
Finite truth tables recover quantifier-free induction. The inherited
scalar construction supplies pairing $\pi$ and projections $L,H$.
For a fixed primitive recursion
\[
 f(0,\vec p)=g(\vec p),\qquad
 f(Sn,\vec p)=h(n,f(n,\vec p),\vec p)
\]
use
\[
\begin{split}
 V(n,\vec p)&=I\bigl(n,\pi(0,g^*(\vec p));
        u.\pi(SL(u),h^*(L(u),H(u),\vec p))\bigr),\\
 f^*(n,\vec p)&=H(V(n,\vec p)).
\end{split}
\]
Equation induction proves the defining equations. This is an external
compiler for fixed descriptions, not an object-language uniform evaluator.

The domain formula is $x=x$ in each direction and equality is unchanged.
Every double-translated operation satisfies its original defining
recurrence. Internal uniqueness identifies the two operations, making
$x=y$ internally total, bijective and operation preserving.

For formulas the logical macros have their ordinary round trips.
The essential paired composite is
\[
 \D_{x,y}B\dashv\vdash
 \exists_*x\bigl(B(x,0)\Rightarrow_*\forall_*y\,B(x,y)\bigr).
\]
Forwards, open D, assume the base, use QE at a fresh parameter and
introduce the macros. Backwards, eliminate the existential and use
the implication and universal instances in DI$_c$.
The Q composite is its conditional universal. Structural induction
supplies both composites for all formulas, including quantified ones.
The explicit falsity bridge preserves contradiction in both directions,
giving equiconsistency.
\end{proof}
This is a written proof using the stated inherited scalar theorems.
It is stronger than mutual interpretability or agreement on standard
numerals, but is not a new proof-assistant verification.

